% Compile: LaTeX
% Biber
% ----------------------------------------------------------------------------
% *** START PACKAGES <<<
% ----------------------------------------------------------------------------
%\documentclass[10pt]{beamer}
\documentclass[pdflatex,compress,8pt,
	xcolor={dvipsnames,svgnames,x11names,table},
	hyperref={colorlinks = true,breaklinks = true, urlcolor = NavyBlue}]{beamer}
\usepackage[utf8]{inputenc}
\usepackage[T2A,T1]{fontenc}
\usepackage[english]{babel}
\usepackage{times}
\usepackage{totcount}
\usepackage{tikz}
	\usetikzlibrary{calc} % для последнего слайда
\regtotcounter{section}
\usepackage{subfig}
\usepackage{fancyhdr}
\usepackage{metalogo}
\usepackage[super]{nth}
\usepackage[font={tiny,it}]{caption} % шрифт подписей к картинкам
	\setbeamertemplate{caption}[numbered]
\usepackage{graphicx}
\usepackage{graphics}
%\usepackage{textcomp}  % numero symbol
\usepackage{gensymb} % degree symbol
%%%%%%%% цветная таблица: начало
\usepackage{booktabs} 
\usepackage{tabularx}
\usepackage{array,etoolbox}
\usepackage{boldline}
\usepackage{colortbl} %% цветная таблица
\setlength{\arrayrulewidth}{0.3mm}
%%%%%% нумерация внутри таблицы %%%
\preto\tabular{\setcounter{magicrownumbers}{0}}
\newcounter{magicrownumbers}
\newcommand\rownumber{\stepcounter{magicrownumbers}\arabic{magicrownumbers}}
\usepackage{longtable}
%%%%%%%%  цветная таблица: конец
\usepackage{wrapfig}
\usepackage{csquotes}

% ------------------------ для разных буллитов в списках -------------------------------
\usepackage{enumitem} 
\usepackage{pifont}
% ----------------------------------------------------------------------------

% Путь к файлам с иллюстрациями
\graphicspath{{figures/}} % path to folder with Figures

% ----------------------------------------------------------------------------
% *** END PACKAGES <<<
% ----------------------------------------------------------------------------
% ----------------------------------------------------------------------------
% *** START SETUP <<<
% ----------------------------------------------------------------------------
% Swiss color Zurich
%\definecolor{color0}{HTML}{150489} % color of title text
%\definecolor{color1}{HTML}{B6CEE5}
%\definecolor{color2}{HTML}{5884B3}
%\definecolor{color3}{HTML}{E5B5C5}
%\definecolor{color4}{HTML}{CC6686}
%\definecolor{color5}{HTML}{F2CEC1}
%\definecolor{color6}{HTML}{E87B70}
%\definecolor{color7}{HTML}{F9EBAA}
%\definecolor{color8}{HTML}{E5CF6C}
%\definecolor{color9}{HTML}{CCE5B5} 
%\definecolor{color10}{HTML}{91BE64}
%\definecolor{color11}{HTML}{B6E3D1}
%\definecolor{color12}{HTML}{5BBE94}
%\definecolor{color13}{HTML}{B6E3D1}
%\definecolor{color14}{HTML}{91BE64}
%\definecolor{color15}{HTML}{CCE5B5}
%\definecolor{color16}{HTML}{E5CF6C}
%\definecolor{color17}{HTML}{F9EBAA}
%\definecolor{color18}{HTML}{E87B70}
%\definecolor{color19}{HTML}{F2CEC1}
%\definecolor{color20}{HTML}{CC6686}
%\definecolor{color21}{HTML}{E5B5C5}
%\definecolor{color22}{HTML}{5884B3}
%\definecolor{color23}{HTML}{B6CEE5}
%\definecolor{color24}{HTML}{5884B3}
%\definecolor{color25}{HTML}{E5B5C5}
%\definecolor{color26}{HTML}{CC6686}
%\definecolor{color27}{HTML}{F2CEC1}

\definecolor{color0}{HTML}{150489} % color of title text
\definecolor{color1}{HTML}{16078A}
\definecolor{color2}{HTML}{240691}
\definecolor{color3}{HTML}{3C049B}
\definecolor{color4}{HTML}{4B03A1}
\definecolor{color5}{HTML}{7100A8}
\definecolor{color6}{HTML}{8004A8}
\definecolor{color7}{HTML}{8D0BA5}
\definecolor{color8}{HTML}{9B169F}
\definecolor{color9}{HTML}{AC2694} 
\definecolor{color10}{HTML}{B6308B}
\definecolor{color11}{HTML}{C33D80}
\definecolor{color12}{HTML}{CE4B75}
\definecolor{color13}{HTML}{D7566C}
\definecolor{color14}{HTML}{DE6164}
\definecolor{color15}{HTML}{E56B5D}
\definecolor{color16}{HTML}{ED7953}
\definecolor{color17}{HTML}{F3864A}
\definecolor{color18}{HTML}{F9973F}
\definecolor{color19}{HTML}{FCA636}
\definecolor{color20}{HTML}{FDB62E}
\definecolor{color21}{HTML}{FDB52E}
\definecolor{color22}{HTML}{FEBB2B}
\definecolor{color23}{HTML}{FCD025}
\definecolor{color24}{HTML}{FBD724}
\definecolor{color25}{HTML}{F8E125}
\definecolor{color26}{HTML}{F4ED27}
\definecolor{color27}{HTML}{F1F426}
\definecolor{color28}{HTML}{F1F426}

\makeatletter

\def\sectioncolor{color0}% color to be applied to section headers

\setbeamercolor{palette primary}{use=structure,fg=structure.fg}
\setbeamercolor{palette secondary}{use=structure,fg=structure.fg!75!black}
\setbeamercolor{palette tertiary}{use=structure,fg=structure.fg!50!black}
\setbeamercolor{palette quaternary}{fg=black}

\setbeamercolor{local structure}{fg=color0}
\setbeamercolor{structure}{fg=color0}
\setbeamercolor{title}{fg=color0}
\setbeamercolor{section in head/foot}{fg=black}

\setbeamercolor{normal text}{fg=black,bg=white}
\setbeamercolor{block title alerted}{fg=red}
\setbeamercolor{block title example}{fg=color0}

% Each \section redefines \sectioncolor and applies the shading with this color
% add as many colors as you need
\AtBeginSection{%
  \renewcommand\sectioncolor{%
    \ifcase\value{section} color0\or color1\or color2\or color3\or color4\or color5\or color6\or color7\or color8\or color9\or color10\or color11\or color12\or color13\or color14\or color15\or color16\or color17\or color18\or color19\or color20\or color21\or color22\else color23\fi}
  \setbeamercolor{frametitle}{fg=red}%
  \pgfdeclareverticalshading{beamer@headfade}{\paperwidth}
  {%
    color(0.25cm)=(bg);
    color(1.25cm)=(\sectioncolor)%
  }
}
\setbeamerfont{frametitle}{size=\normalsize}

\newlength\sectionboxwd

\AtBeginDocument{%
  \ifnum\totvalue{section}>0 
    \setlength\sectionboxwd{\dimexpr\paperwidth/\totvalue{section}\relax}
  \else
    \setlength\sectionboxwd{\paperwidth}
  \fi
} 

\newcommand\insertcolors{%
  \ifnum\totvalue{section}>0
    \setlength\fboxsep{0pt}%
    \foreach \x in {1,...,\totvalue{section}}
    {\colorbox{color\x}{\phantom{\rule{\dimexpr\the\sectionboxwd\relax}{5ex}}}}%
  \else\fi
}

\setbeamertemplate{section in head/foot}{\setlength\fboxsep{0pt}\hspace{-1.875ex}%
    \parbox[c][4ex][t]{\sectionboxwd}{%
      \hfill\parbox{\dimexpr\sectionboxwd-8pt\relax}{%
      \raggedright\insertsectionhead%
      }\hfill\mbox{}%
    }%
}

\setbeamertemplate{headline}
{%
  \begin{beamercolorbox}[wd=\paperwidth]{section in head/foot}
    \insertcolors%
    \vskip-4ex%
    \insertsectionnavigationhorizontal{\paperwidth}{\hskip-1.875ex}{\hskip-8pt}\vskip2pt
  \end{beamercolorbox}%
}


\pgfdeclareverticalshading{beamer@headfade}{\paperwidth}
{% полоска вверху
  color(0.25cm)=(bg);
  color(1.25cm)=(color18)% титульный лист цвет
}

\addtoheadtemplate{}{\pgfuseshading{beamer@headfade}\vskip-1.25cm}

\setbeamertemplate{navigation symbols}{}
\addtobeamertemplate{footline}{}{%
  \hfill\color{\sectioncolor}\rule[0.5cm]{2.5cm}{1pt}\rule[0.5cm]{1pt}{1cm}\hspace*{0.5cm}}
  
 \setbeamerfont{institute}{size=\normalfont}
% ----------------------------------------------------------------------------
% *** END SETUP <<<
% ----------------------------------------------------------------------------

% ----------------------------------------------------------------------------
% *** START BIBLIOGRAPHY <<<
% ----------------------------------------------------------------------------
\usepackage[
	backend=biber, 
	style = numeric,
	maxbibnames=99,
	citestyle=numeric,
	giveninits=true,
	isbn=true,
	url=true,
	natbib=true,
	sorting=ndymdt,
	bibencoding=utf8,
	useprefix=false,
	language=auto, 
	autolang=other,
	backref=true,
	backrefstyle=none,
	indexing=cite,
]{biblatex}
\DeclareSortingTemplate{ndymdt}{
  \sort{
    \field{presort}
  }
  \sort[final]{
    \field{sortkey}
  }
  \sort{
    \field{sortname}
    \field{author}
    \field{editor}
    \field{translator}
    \field{sorttitle}
    \field{title}
  }
  \sort[direction=descending]{
    \field{sortyear}
    \field{year}
    \literal{9999}
  }
  \sort[direction=descending]{
    \field[padside=left,padwidth=2,padchar=0]{month}
    \literal{99}
  }
  \sort[direction=descending]{
    \field[padside=left,padwidth=2,padchar=0]{day}
    \literal{99}
  }
  \sort{
    \field{sorttitle}
  }
  \sort[direction=descending]{
    \field[padside=left,padwidth=4,padchar=0]{volume}
    \literal{9999}
  }
}

\addbibresource{AGUBib.bib}
%\renewcommand*{\bibfont}{\footnotesize}
\renewcommand*{\bibfont}{\scriptsize}

%%%%%%%%%%%% значок Open Access: начало %%%%%%%
%\usepackage[tikzsymbol=plos]{biblatex-ext-oa}
%\DeclareOpenAccessEprintUrl[always]{hal}{%
 % http://hal.archives-ouvertes.fr/\thefield{eprint}}
%\DeclareOpenAccessEprintAlias{HAL}{hal}
%%%%%%%%%%% значок Open Access:: конец %%%%%%%

\setbeamertemplate{bibliography item}{\insertbiblabel}
%\setbeamertemplate{bibliography item}{}

% ----------------------------------------------------------------------------
% *** END BIBLIOGRAPHY <<<
% ----------------------------------------------------------------------------
%%%%%%%%% hyper ref setup %%%%%%%%%
\usepackage[colorlinks=true]{hyperref}
\hypersetup{pdftitle={Accurate and rapid big spatial data processing by scripting cartographic algorithms:
advanced seafloor mapping of the deep-sea trenches along the margins of the Pacific Ocean}, 
	pdfauthor={Lemenkova Polina}, 
	pdfsubject={PhD Thesis}, 
	pdfcreator={Lemenkova Polina}, 
	pdfproducer={Lemenkova Polina}, 
	colorlinks=true,linkcolor=blue, 
	citecolor=NavyBlue, 
	filecolor=magenta, 
	urlcolor=blue,
	pdfkeywords={Pacific Ocean, hadal trench, deep-sea trench, geomorphology, bathymetry, cartography, mapping, GMT, Generic Mapping Tools, Python, R, programming language, Matplotlib, Seaborn, submarine geomorphology, data analysis, statistics, spatial analysis, comparative analysis, statistical analysis, cluster analysis, factor analysis, algorithm, oceanography, GIS, marine geology, geomorphology}
	}	

%%%%%%%% цветные разделители
\usepackage[object=vectorian]{pgfornament}
%%%%%%%% цветные разделители

%%%%%%%%%%%%%%%%%%%%%%%%%%%%%%%%%%%
% ----------------------------------------------------------------------------
% *** TITLE <<<
% ---------------------------------------------------------------------------
\vspace{1em}
\title{Accurate and rapid big spatial data processing \\
by scripting cartographic algorithms: \\
advanced seafloor mapping of the deep-sea trenches \\along the margins of the Pacific Ocean}
\subtitle{Presented at \\
\href{https://agu.confex.com/agu/21workshop2/meetingapp.cgi/Home/0}{Bringing Land, Ocean, Atmosphere and Ionosphere Data \\to the Community for Hazard Alerts}\\
American Geophysical Union (AGU) \\
Online: 24-28 May 2021 \\
Section 'Big Data and Machine Learning'}
\author{Polina Lemenkova}
%\institute{Schmidt Institute of Physics of the Earth, Russian Academy of Science (IPE RAS)\\
%\small{Department of Natural Disasters, Anthropogenic Hazards and Seismicity of the Earth. \\
%Laboratory of Regional Geophysics and Natural Disasters (Nr. 303)}}
\date{May 24, 2021}
%%%%%%%%%%%%%%%%%%%%%%%%%%%%%%%%%%%
\logo{\includegraphics[width=.25\textwidth]{AGU100_logo_H-CMYK.png}\vspace{-35pt}\hspace{10pt}}
%%%%%%%%%


%%%%%%%%%%%%%%%%%%%%%%%%%%%%%%%%%%%
\begin{document}

% ----------------------------------------------------------------------------
% *** Titlepage <<<
% ----------------------------------------------------------------------------
\maketitle
% ----------------------------------------------------------------------------
% *** END of Titlepage >>>
% ----------------------------------------------------------------------------
\setbeamertemplate{footline}[text line]{%
\parbox{\linewidth}{\vspace*{-8pt}Polina Lemenkova: \emph{Accurate and rapid big spatial data processing by scripting cartographic algorithms...}, AGU, 24/05/2021} \hfill\insertpagenumber}
%\setbeamertemplate{headline}{\insertpagenumber}


% ----------------------------------------------------------------------------
% *** END of Footnote >>>
% ----------------------------------------------------------------------------
% ----------------------------------------------------------------------------
% *** Introduction <<<
% ----------------------------------------------------------------------------
\section{Int}

\begin{frame}\frametitle{Introduction}
\begin{exampleblock}{What is Big Earth Data ?}
\begin{itemize}
        	\item \textcolor{red}{Concept}: Large and complex massifs (collections) of digital geospatia data
	\item \textcolor{red}{Features}: massive, multi-source, heterogeneous, multi-temporal, multi-scalar, highly dimensional, highly complex, nonstationary, and unstructured \cite{Guo2017}
	\item \textcolor{red}{Nature}: Big Earth data are associated with geosciences (geology, geography)
	\item \textcolor{red}{VVV}: Three major magic 'V' characterising Big Data: 
		\begin{itemize}\small
      		  	\item \textcolor{ProcessBlue}{volume} (unprecedented wealth of data)
			\item \textcolor{ProcessBlue}{velocity} (growth rate and the demand for high-speed processing and mapping)
			\item \textcolor{ProcessBlue}{variety} (multi-source heterogeneous spatial data
		\end{itemize}
	\item \textcolor{red}{Complexity}: need for simultaneous handling of various types of structured and non-structured data \cite{Yangetal2019}
	\item \textcolor{red}{Origin}: Big Earth data are derived from Earth observations (EO) captured by Remote Sensing (RS) methods \cite{Guoetal2016}
	\item \textcolor{red}{Sources}: from orbiting satellites (Landsat TM/ETM+, Sentinel-2A, SPOT, etc), weather monitoring stations, ecological observatories, seismic stations, model simulations and forecasts, atmospheric and oceanic gridded data, time series observations and other data streams \cite{Sellarsetal2013}.
	\item \textcolor{red}{Importance}: The value of big Earth data processing to geoscience research: development of \emph{data-driven science} \cite{Lemenkova202101}
\end{itemize}
\end{exampleblock}

\end{frame}

% ---------------------------------------------------------------------------- 
\begin{frame}\frametitle{Introduction}

\begin{exampleblock}{Big Spatial Data as an Opportunity}
\begin{itemize}	
	\item \textcolor{red}{Possibilities}: Technical revolution in the amount of Earth data (satellite imagery, digital raster and vector datasets, tabular data) enabled smooth and quick access to monitoring of the large regions of the Earth, such as Pacific Ocean \cite{Lemenkova202106}
	\item \textcolor{red}{Access}: Open datasets and repositories available online enable to access, download, examine, and map big spatial data for various \\geospatial applications \cite{Baumannetal2015}, \cite{Lemenkova202103}
	\item \textcolor{red}{Coverage}: Big Earth data provide information on spatial events across local, regional, and global scales 
\end{itemize}
\end{exampleblock}

\begin{exampleblock}{Machine Learning for Processing big Earth data}
\begin{itemize}
	\item \textcolor{blue}{Problem}: Extracting knowledge and information from big Earth data requires special technical tools $=>$ need for Machine Learning (ML) algorithms and approaches to data handling \cite{Hanetal2019}
	\item \textcolor{red}{Paradigm}: ML is a method of Earth data analysis that automates data processing, modeling, mapping and visualization \cite{Lemenkova202021}, \cite{Lemenkova202102}. 
	\item \textcolor{red}{Hierarchy}: ML is a branch of Artificial Intelligence (AI). It is based on the assumption that GIS can \emph{learn from data and extract information}, identify geospatial patterns and make decisions (e.g. risk assessment, tsunami, earthquakes, landslides)
\end{itemize}
\end{exampleblock}
	
\end{frame}

% ----------------------------------------------------------------------------
\begin{frame}\frametitle{Research Highlights}
 
\begin{alertblock}{Machine Learning in big Earth data processing}
\begin{itemize}
	\item \textcolor{red}{Core Idea}: \emph{minimal} human intervention, \emph{maximal} machine-based processing
	\item \textcolor{blue}{Examples}: programming languages (Python, Julia, R, Octave), scripting GIS and toolsets (GRASS GIS, GMT)
	\item \textcolor{red}{Applications}: Geotectonics, geology, geophysics, sedimentation, seafloow mapping, 3D modeling, geodynamics, environment  \cite{Amanietal2019}, \cite{Lemenkova202104}, \cite{Lemenkova202105} for automatization of big Earth data processing
	\item \textcolor{red}{Case}: This research uses Generic Mapping Tools (GMT) scripting toolset that is based on special syntax enabling to quickly operate with big spatial data using shell scripts from the console \cite{Lemenkova202028}, \cite{Lemenkova202025}, \cite{Lemenkova202020}.
\end{itemize}
\end{alertblock}

\begin{exampleblock}{Challenge of the Research Topic}
\begin{itemize}
        	\item Using scripting algorithms for mapping deep-sea trenches of the Pacific Ocean
\end{itemize}
\end{exampleblock}
	
\begin{alertblock}{Research Techniques}
\begin{itemize}
	\item This research includes algorithms of Generic Mapping Tools (GMT) that has a scripting syntax enabling to quickly operate with spatial data using shell scripts
\end{itemize}
\end{alertblock}
\end{frame}

% ----------------------------------------------------------------------------
\section{PO}

\begin{frame}\frametitle{Pacific Ocean}
\vspace{1em}
\begin{figure}[H]
	\centering
		\includegraphics[width=8cm]{Fig-PO.jpg}
	\caption{General bathymetric map of the seafloor of the Pacific Ocean, ETOPO1}\label{fig:PO.jpg}
\end{figure}\small
The \textcolor{red}{biggest} ocean on the Earth, Pacific Ocean differs in \textcolor{DarkOrchid}{western} and \textcolor{DarkOrchid}{eastern} parts according to relief, geotectonics and geological evolution.
Western Pacific (WP) and Eastern Pacific (EP) are different basins, separated by the East Pacific Rise. 
\end{frame}

% ----------------------------------------------------------------------------
\begin{frame}\frametitle{Pacific Ocean}
\vspace{1em}
\begin{figure}[H]
	\centering
		\includegraphics[width=8.0cm]{Fig-POg.jpg}
	\caption{General geologic map of the seafloor of the Pacific Ocean}\label{fig:POg.jpg}
\end{figure}
\small{Pacific Ocean is known for variability of morphostructures, complexity of tectonic features, high seismicity and magmatism, lithology and geologic evolution. \\Pacific Ocean differs from others: enormous size (ca. 40\% of the Earth's surface and ca. 50\% of the World Ocean, 178.7 M $km^2$), geologic structural complexity and bathymetric asymmetry}
\end{frame}

% ----------------------------------------------------------------------------
\begin{frame}\frametitle{Pacific Ocean}
\begin{alertblock}{Research Highlights: Trenches of the Pacific Ocean}
\begin{itemize}\footnotesize
            \item Trenches of the Pacific Ocean are located along its margins and have diverse geomorphology. 
            \item Geomorphology of trenches reflects a variety of impact factors: tectonic, geographic and geologic setting
\end{itemize}
\end{alertblock}	

\begin{table}\scriptsize
\caption{General statistics on major 20 deep-sea trenches of the Pacific Ocean}
\begin{tabular}{@{\makebox[2em][r]{\rownumber\space}} | p{3.5cm} p{1.2cm} p{1.2cm} p{0.7cm} |}
\multicolumn{1}{@{\makebox[2em][r]{No}} | l}{Name} & Max depth & Length* & Width* \\ 
		\arrayrulecolor[RGB]{184,210,0}\hline\hline
	Aleutian & 8,109 & 3,400 & 59 \\
       		\arrayrulecolor[RGB]{195,216,37}\hline
	Mariana & 11,034 & 2,550 & 59 \\
        		\arrayrulecolor[RGB]{170,207,83}\hline
	Philippine & 10,540 & 1,320 & 65 \\
        		\arrayrulecolor[RGB]{199,220,104}\hline
	Kuril-Kamchatka & 10,542 & 2,900 & 59 \\
        		\arrayrulecolor[RGB]{216,230,152}\hline
	Middle America & 6,669 & 2,750 & 34 \\
        		\arrayrulecolor[RGB]{224,235,175}\hline
	Peru-Chile & 8,065 & 5,900 & 64 \\
        		\arrayrulecolor[RGB]{185,208,139}\hline	
	Hikurangi & 3,750 & 350 & 80 \\
		\arrayrulecolor[RGB]{185,208,139}\hline	
	Puysegur & 6,300 & 620 & 76 \\
        		\arrayrulecolor[RGB]{185,208,139}\hline	
	Palau & 10,994 & 700 & 47 \\
        		\arrayrulecolor[RGB]{185,208,139}\hline
	Japan & 8,513 & 800 & 59 \\
        		\arrayrulecolor[RGB]{185,208,139}\hline
	Kermadec & 10,047 & 1,200 & 88 \\
        		\arrayrulecolor[RGB]{185,208,139}\hline
	Tonga & 10,882 & 1,375 & 78 \\
        		\arrayrulecolor[RGB]{185,208,139}\hline
	Izu-Bonin & 9,780 & 2,800 & 82 \\
        		\arrayrulecolor[RGB]{185,208,139}\hline
	New Britain (Bougainville) & 9,140 & 335 & 70 \\
        		\arrayrulecolor[RGB]{185,208,139}\hline
	San Cristobal & 8,255 & 742 & 64 \\
        		\arrayrulecolor[RGB]{185,208,139}\hline
	Manila & 5,400 & 648 & 76 \\
        		\arrayrulecolor[RGB]{185,208,139}\hline
	Yap & 8,850 & 500 & 45 \\
        		\arrayrulecolor[RGB]{185,208,139}\hline
	Vanuatu (New Hebrides) & 7,600 & 1,200 & 34 \\
        		\arrayrulecolor[RGB]{185,208,139}\hline
	Vityaz & 5,984 & 580 & 65 \\
        		\arrayrulecolor[RGB]{185,208,139}\hline
	Ryukyu & 7,460 & 2,250 & 38 \\
        		\arrayrulecolor[RGB]{131,155,92}\hline\hline
\end{tabular}
\end{table}
*\footnotesize{Length and width are indicated in km}
\end{frame}

% ----------------------------------------------------------------------------
\section{Data}

\begin{frame}\frametitle{Data}
\vspace{1em}
\begin{itemize}\footnotesize
            \item Understanding complexity and variability of trenches requires processing of large amounts of diverse datasets 
            \item Datasets used in this project are non-coherent and multi-source 
            \item Data should be handled effectively, rapidly yet accurately.
            \item High-resolution datasets aimed at print-quality mapping and modelling. Accuracy of digital Earth data is of crucial importance \cite{Smith}, because it directly controls research output. 
\end{itemize}

\begin{table}\footnotesize\caption{Data used in this project: their sources, types and precision}\label{tab:T2}
\begin{tabular}{@{\makebox[2em][r]{\rownumber\space}} | p{1.5cm} p{1.5cm} p{5cm} }
\multicolumn{1}{@{\makebox[2em][r]{No}} | l}{Data} & Origine & Type, resolution \\ 
		\arrayrulecolor[RGB]{211,204,214}\hline\hline		
        ETOPO1 & NOAA & 1 arc-min GRM grid \cite{AmanteEakins2009}\\       
       		\arrayrulecolor[RGB]{166,154,189}\hline
	ETOPO5 & NOAA & 5 arc min GRM grid \\       
       		\arrayrulecolor[RGB]{166,154,189}\hline
	GEBCO & BODC & 15 arc sec DEM raster grid \\       
       		\arrayrulecolor[RGB]{166,154,189}\hline				
        SRTM & NASA & 15-sec DEM raster grid \cite{Beckeretal2009} \\      
        		\arrayrulecolor[RGB]{166,165,196}\hline\hline
	Geology & USGS & Vector layers \\      
        		\arrayrulecolor[RGB]{166,165,196}\hline\hline
	Gravity & Scripps IO & CryoSat-2, Jason-1 grid \cite{Sandwelletal2013}\\   
		\arrayrulecolor[RGB]{166,165,196}\hline\hline
	EGM96 & Scripps IO & Geopotential geoid model \cite{Lemoineetal1998} \\      
        		\arrayrulecolor[RGB]{166,165,196}\hline\hline
	EGM-2008 & Scripps IO & Geopotential geoid model \cite{Pavlisetal2012} \\      
        		\arrayrulecolor[RGB]{166,165,196}\hline\hline	
	GlobSed & NOAA & Total Sediment Thickness of the World's Oceans and Marginal Seas  \\      
        		\arrayrulecolor[RGB]{166,165,196}\hline\hline	
	GLOBE dataset & NOAA & Topography model \\      
        		\arrayrulecolor[RGB]{166,165,196}\hline\hline		
	ASCII & Scripps IO & Topo tables (xyz format) \\      
        		\arrayrulecolor[RGB]{166,165,196}\hline\hline
	GEBCO & IHO-IOC & Geographic gazetteer \cite{IHOIOC2012} \\      
        		\arrayrulecolor[RGB]{166,165,196}\hline\hline
\end{tabular}
\end{table}
\footnotesize{*30 arc second = ca. 1km cell size, 15 arc sec = ca. 500 m cell size, etc.}
\end{frame}

% ----------------------------------------------------------------------------
\section{DR}

\begin{frame}\frametitle{GEBCO}
\begin{minipage}[0.4\textheight]{\textwidth}
\begin{columns}[T]
\begin{column}{0.4\textwidth}
\vspace{2em}
\begin{dinglist}{51}\small
     \item Satellite Altimetry and Earth Observation (EO) have generated huge massives of data stored in big data repositories
     \item GEBCO produces and makes open a range of bathymetric datasets (GEBCO-2020). 
     \item GEBCO global gridded bathymetric datasets; GEBCO Gazetteer of Undersea Feature Names; GEBCO world map 
     \item Advanced cartographic methodologies of GMT are capable of handling the size (e.g. large GEBCO grid), complexity of formats (with help of reformatting by AWK) and variety of data, as a response to such needs. 
     \item Data manipulation, mapping, information extraction (metadata), visualisation $=>$ processes of big Earth data handling.            
\end{dinglist}	
\end{column}
\begin{column}{0.6\textwidth}
\vspace{2ex}
\begin{figure}[H]
	\centering
		\includegraphics[width=6.5cm]{Fig-KKT1.jpg}
	\caption{Example of the GEBCO grid used for mapping trench}\label{fig:KKT1}
\end{figure}
\end{column}
\end{columns}
\end{minipage}
\end{frame}

% ----------------------------------------------------------------------------
\section{ML}
\begin{frame}\frametitle{Octave/MATLAB}
\vspace{2em}
\begin{minipage}[0.4\textheight]{\textwidth}
\begin{columns}[T]
\begin{column}{0.5\textwidth}
\begin{figure}[H]
	\centering
		\includegraphics[width=6cm]{S02.jpg}
	\caption{Snippet of the Octave/MATLAB script used for modeling the Kuril-Kamchatka Trench}\label{fig:S2}
\end{figure}	
\end{column}

\begin{column}{0.4\textwidth}
\begin{figure}[H]
	\centering
		\includegraphics[width=5cm]{Fig-KKT6.jpg}
	\caption{Cross-section profiles of the Kuril-Kamchatka Trench. Plotting: Octave}\label{fig:KKT6}
\end{figure}
\begin{dinglist}{224}\scriptsize
     \item Modelling geospatial data using ML methods and programming languages enables to detect hidden correlations and patterns that a human may not see. 
     \item Inaccessible location $=>$ seafloor of the deep-sea trenches can only be visualized using modelling tools and programming approach
\end{dinglist}
\end{column}
\end{columns}
\end{minipage}
\end{frame}

% ----------------------------------------------------------------------------

\section{AT}
\begin{frame}\frametitle{Aleutian Trench}
\vspace{1em}
\begin{figure}[H]
	\centering
		\includegraphics[width=7cm]{Fig-A1.jpg}
	\caption{ETOPO1 grid for topographic mapping: example of Aleutian Trench}\label{fig:AT1}
\end{figure}
Solutions of GMT: 
 \begin{dingautolist}{202}\footnotesize
     \item Advanced level of cartographic functionality: variety of projections, data processing, import/export, etc
     \item High-quality cartographic solutions: colour palettes, layouts, grids, vector lines, advanced design approaches
     \item Flexibility of coding / shell scripting from UNIX console
     \item Embedded data analysis and visualisation: mapping, modelling, logical queries, \\ statistical analysis, import/export, vector and raster processing
\end{dingautolist}
\end{frame}

% ----------------------------------------------------------------------------
\begin{frame}\frametitle{Data Integration}
\vspace{1em}
Multi-source data mixing and integration using GMT adjustments:
 \begin{dingautolist}{172}\small
     \item Open Data Mining (e.g. GEBCO, ETOPO1, geological datasets and other open sources) 
     \item Predictive analytics in geology: forecasting and prognosis in hazard risk assessment
     \item Simulation modeling in Earth sciences: earthquakes, landslides, volcanic activity, seismicity
     \item Data analysis by GMT has potential to advance development of cartography
     \item It enables scientific discoveries, based on spatial analysis in geology
\end{dingautolist}	
   
\begin{figure}[H]
	\centering
		\includegraphics[width=6cm]{Fig-MAT5.jpg}
	\caption{Map showing geological settings in the Middle America Trench}\label{fig:MAT5}
\end{figure}
\end{frame}

% ----------------------------------------------------------------------------
\begin{frame}\frametitle{Middle America Trench}
\vspace{2em}
\begin{minipage}[0.4\textheight]{\textwidth}
\begin{columns}[T]
\begin{column}{0.4\textwidth}
\begin{figure}[H]
	\centering
		\includegraphics[width=6cm]{Fig-MAT10.jpg}
	\caption{Cross-section profiles of the Middle America Trench (Guatemala segment)}\label{fig:MAT10}
\end{figure}
\end{column}

\begin{column}{0.4\textwidth}
\vspace{2em}
An example of data analysis \\using GMT: Guatemala Trench: 
 \begin{dinglist}{52} \scriptsize
     \item MAT geometry demonstrates its relatively steep and strait shape. 
     \item major depth samples: between the -3000 and -6200 m; seafloor 3-5 km wide.
     \item frequency samples distribution shows changes by profiles. Increase of frequency between -3000 to -4000 m caused by the geologic settings of the Cocos Plate. 
     \item the deepest depth data range with 1,111 records takes 25\% of the samples. 
     \item mean sample values: -2,200, and 7,5\% at -4,000 to -3,500 m
     \item median values: -3,200 to -3,300 m. 
 \end{dinglist}  		
\end{column}
\end{columns}
\end{minipage}
\end{frame}

% ----------------------------------------------------------------------------
\section{MA}
\begin{frame}\frametitle{Middle America Trench}
\vspace{2em}
\begin{minipage}[0.4\textheight]{\textwidth}
\begin{columns}[T]
\begin{column}{0.4\textwidth}
\begin{figure}[H]
	\centering
		\includegraphics[width=6cm]{S01.jpg}
	\caption{Snippet of the GMT script for mapping Middle America Trench}\label{fig:S01}
\end{figure}
\end{column}

\begin{column}{0.5\textwidth}
\begin{figure}[H]
	\centering
		\includegraphics[width=6cm]{Fig-MAT1.jpg}
	\caption{Topographic map and location of the Middle America Trench}\label{fig:MAT1}
\end{figure}	
\begin{itemize}\footnotesize
            \item Processing large volumes of heterogeneous and rapidly arriving digital geoinformation using traditional GIS is time-consuming. 
            \item GMT-based mapping uses scripts which optimises cartographic workflow
\end{itemize}
\end{column}
\end{columns}
\end{minipage}
\end{frame}

% ----------------------------------------------------------------------------
\section{PC}

\begin{frame}\frametitle{Peru-Chile Trench}
\vspace{2em}
\begin{minipage}[0.4\textheight]{\textwidth}
\begin{columns}[T]
\begin{column}{0.4\textwidth}
\begin{figure}[H]
	\centering
		\includegraphics[width=3.8cm]{Fig-PCT1.jpg}
	\caption{Topographic map of the PCT}\label{fig:PCT1}
\end{figure}
\end{column}

\begin{column}{0.6\textwidth}

\begin{alertblock}{GMT: advantages in big Earth data processing}
\begin{itemize}\footnotesize
	\item \textcolor{red}{Data structure}: Structured big data are stored in the relevant folders known by GMT
	\item \textcolor{red}{Data management}: Earth data are managed by specific GMT modules that support visualisation
	\item \textcolor{red}{Flexibility}: certain cartographic elements can be conveniently extracted for only the part of the data relevant to study area
\end{itemize}
\end{alertblock}

\begin{exampleblock}{Big Earth data technologies benefit geoscience research}
\begin{dinglist}{43}\footnotesize
     \item Visualisation of the analytical data: e.g. bathymetric analysis
     \item Analysis of big Earth data reveals the topography of the Peru-Chile Trench (PCT), a.k.a. Atacama Trench
     \item PCT - the longest Pacific trench (5900 km)
     \item PCT stretches from Ecuador to Chile and covers area of ca. 590,000 km$^2$
     \item PCT is narrower than other Pacific trenches (64 km). 
     \item Max depth 8065 m $=>$ deepest trench in South PO        
\end{dinglist}
\end{exampleblock}	

\end{column}
\end{columns}
\end{minipage}
\end{frame}

% ----------------------------------------------------------------------------
\begin{frame}\frametitle{Peru-Chile Trench}
\vspace{2em}
\begin{minipage}[0.4\textheight]{\textwidth}
\begin{columns}[T]
\begin{column}{0.4\textwidth}
\begin{figure}[H]
	\centering
		\includegraphics[width=5.5cm]{S03.jpg}
	\caption{Snippet of the GMT script for mapping Peru-Chile Trench}\label{fig:S03}
\end{figure}
\end{column}
\begin{column}{0.5\textwidth}
\begin{figure}[H]
	\centering
		\includegraphics[width=6cm]{Fig-PCT6.jpg}
	\caption{Cross-section profiles of the Peru-Chile Trench}\label{fig:PCT6}
\end{figure}
\begin{dingautolist}{192}\scriptsize
	\item Automatic modelling of trenches to study landforms, detect correlation with geologic and tectonic settings
	\item Automatization in data analysis: automated digitizing of the cross-section profiles
\end{dingautolist}
\end{column}
\end{columns}
\end{minipage}
\end{frame}

 % ----------------------------------------------------------------------------
\section{H}
\begin{frame}\frametitle{Hikurangi Trench}
\vspace{2em}
\begin{minipage}[0.4\textheight]{\textwidth}
\begin{columns}[T]
\begin{column}{0.4\textwidth}
\begin{figure}[H]
	\centering
		\includegraphics[width=5.0cm]{Fig-H1.jpg}
	\caption{Topographic settings of the Hikurangi Trench.}\label{fig:H1}
\end{figure}
\end{column}
\begin{column}{0.6\textwidth}
\begin{figure}[H]
	\centering
		\includegraphics[width=6.0cm]{Fig-H2.jpg}
	\caption{Modeled cross-section profiles of the Hikurangi Trench}\label{fig:H2}
\end{figure}
\begin{dinglist}{105}\scriptsize
     \item Various cartographic techniques can be used for modelling and mapping of deep-sea trenches in the Pacific Ocean.
     \item Advanced cartographic techniques consist in two important characteristics: they are \textcolor{red}{script-based} and \textcolor{red}{data-driven}
     \item Cartographic scripting is a console-based technique that provides detailed stepwise controlling of the cartographic workflow enabling rapid mapping    
     \item Rapid mapping by GMT enables visualisation and representation of the tectonic processes in the seafloor 
     \item Accurate mapping helps more thoughtful interpretation of the underlying tectonic processes affecting the geomorphology of the trench. 
     \item Data-driven techniques use Earth big data to capture the correlation between geological, tectonic and bathymetric variables
\end{dinglist}
\end{column}
\end{columns}
\end{minipage}
\end{frame}

% ----------------------------------------------------------------------------
\section{Ps}

\begin{frame}\frametitle{Puysegur Trench}
\vspace{2em}
\begin{minipage}[0.4\textheight]{\textwidth}
\begin{columns}[T]
\begin{column}{0.4\textwidth}
\begin{figure}[H]
	\centering
		\includegraphics[width=4.5cm]{Fig-P1.jpg}
		\caption{Geologic map: Puysegur and Hjort trenches}\label{fig:P1}
\end{figure}
\end{column}

\begin{column}{0.6\textwidth}
\begin{figure}[H]
	\centering
		\includegraphics[width=5.5cm]{Fig-P2.jpg}
		\caption{Modeled cross-section profiles of the Puysegur and Hjort trenches}\label{fig:P2}
\end{figure}
 \begin{dinglist}{234}\scriptsize
     \item GMT-based mapping techniques have the advantage of increased speed, accuracy and precision of mapping which enables rapid big data processing
     \item Validity and utility of data models in GMT are based on embedded mathematical algorithms of data modelling and cartographic principles enabling deeper understanding of the underlying processes affecting the shape of the trench
\end{dinglist}	
\end{column}
\end{columns}
\end{minipage}
\end{frame}

 % ----------------------------------------------------------------------------
\section{K}

\begin{frame}\frametitle{Kermadec Trench}
\vspace{2em}
\begin{minipage}[0.4\textheight]{\textwidth}
\begin{columns}[T]
\begin{column}{0.4\textwidth}
\begin{figure}[H]
	\centering
		\includegraphics[width=6cm]{Fig-TK1.jpg}
	\caption{Topographic map and location of the Kermadec and Tonga trenches. \cite{Lemenkova201994}}\label{fig:TK1}
\end{figure}
\end{column}
\begin{column}{0.38\textwidth}

 \begin{dinglist}{46}\footnotesize
     \item In order to analyse the geomorphology of the seafloor landforms, it is important to use  \textcolor{red}{powerful cartographic tools}, such as  \textcolor{red}{GMT} that can handle \textcolor{red}{big Earth data}
      \item \textcolor{red}{Modelling links} among seafloor geomorphology, tectonics and geology requires more deep insights into geologic evolution of the Pacific Ocean and current topography of the seabed
      \item However, visualisation of the hidden relief of the seafloor can be difficult without \textcolor{red}{high-tech cartographic solutions}
      \item Cartographic applications of machine learning as scripts provides better techniques in understanding the links of geology, seismicity and bathymetry 
      \item \textcolor{red}{Accurate and rapid mapping} enables to efficiently analyse impact factors controlling geomorphological structures through observing big Earth data related to interactions among seafloor geomorphology and tectonics 
 \end{dinglist}	
\end{column}
\end{columns}
\end{minipage}
\end{frame}


 % ----------------------------------------------------------------------------
\section{T}

\section{TT}
\begin{frame}\frametitle{Tonga Trench}
\vspace{2em}
\begin{minipage}[0.4\textheight]{\textwidth}
\begin{columns}[T]
\begin{column}{0.4\textwidth}
\begin{figure}[H]
	\centering
		\includegraphics[width=6cm]{Fig-TK3.jpg}
	\caption{Cross-section profiles of the Tonga and Kermadec trenches. \cite{Lemenkova201994}}\label{fig:TK3}
\end{figure}
\end{column}
\begin{column}{0.4\textwidth}

\begin{dinglist}{47}\footnotesize
     \item Applying machine learning and scripting techniques of GMT in cartography for mapping seafloor helps tackling problems relating to \textcolor{red}{analysing, mapping, visualising and cross-section modelling} of the submarine geomorphology of the trenches.
      \item Big Earth data come from different \textcolor{red}{heterogeneous sources} and are spatially not uniform. 
      \item Big Earth data are of \textcolor{red}{varying spatio-temporal resolution and origin} which understandably poses some problems to use raw data for mapping and requires \textcolor{red}{data pre-processing, sorting and management}.
      \item Therefore, \textcolor{red}{common resolution} is required in order to perform Earth \textcolor{red}{data integration} and spatial analysis in a project of Pacific Ocean 
\end{dinglist}  	

\end{column}
\end{columns}
\end{minipage}
\end{frame}

 % ----------------------------------------------------------------------------
\section{Vn}

\begin{frame}\frametitle{Vanuatu Trench}
\begin{minipage}[0.4\textheight]{\textwidth}
\begin{columns}[T]
\begin{column}{0.5\textwidth}
\vspace{2em}
\begin{figure}[H]
	\centering
	\includegraphics[width=6.0cm]{Fig-VV1.jpg}\caption{Topographic map and location of the Vanuatu (New Hebrides) and Vityaz trenches}
	\label{fig:VV1}
\end{figure}
\end{column}
\begin{column}{0.5\textwidth}
\vspace{2em}
Key points on the formation of the Vanuatu (a.k.a. New Hebrides) Trench: 
\begin{itemize}
            \item located in Fiji Basin, East Australia
            \item area of the double convergent complex boundary 
            \item PP subducts under the IAP westwards
            \item IAP subducts beneath North Fiji back-arc basin northeastwards. 
            \item Complex geophysical settings, high seismicity, geodynamic complexity and instability of the region 
            \item Vanuatu arc stretches ca. 1200 km in SW PO, 10\degree S-23\degree S, 160\degree -180\degree E. 
            \item Vanuatu trench-arc system divides the NH back-arc troughs and active North Fiji marginal basin.
\end{itemize}
\end{column}
\end{columns}
\end{minipage}
\end{frame}

\begin{frame}\frametitle{Vanuatu Trench}
\begin{minipage}[0.4\textheight]{\textwidth}
\begin{columns}[T]
\begin{column}{0.5\textwidth}
\vspace{2em}
\begin{figure}[H]
	\centering
		\includegraphics[width=6cm]{Fig-VV9.jpg}
	\caption{Modeled cross-sections of the Vanuatu (New Hebrides) and Vityaz trenches}\label{fig:VV9}
\end{figure}
\end{column}
\begin{column}{0.5\textwidth}
\vspace{2em}
Key points on the comparative analysis of the geomoprhic modelling of Vityaz and Vanuatu trenches: 
\begin{itemize}
            \item Vt has more flat wide bottom similar to a trough, with steeper gradient slope on the eastern flank. 
            \item Vn has more V-form classic shape for the trench with gentle shapes on both western and eastern slopes.
            \item Vt has shallower depths with maximal not exceeding -5,000 m while Vn is more deep with -6,000 m. 
            \item The surrounding relief of the Vn \& Vt varies: Vt has the flat neighboring surface; 
            \item Vn is surrounded by complex submarine relief of North Fiji Basin (its eastern flank) and its western flank
\end{itemize}
\end{column}
\end{columns}
\end{minipage}
\end{frame}

\section{Vt}

\begin{frame}\frametitle{Vityaz Trench}
\begin{minipage}[0.4\textheight]{\textwidth}
\begin{columns}[T]
\begin{column}{0.5\textwidth}
\vspace{2em}
\begin{figure}[H]
	\centering
		\includegraphics[width=6cm]{Fig-VV2.jpg}
	\caption{Geologic map and tectonic setting of the Vanuatu (New Hebrides) and Vityaz trenches}\label{fig:VV2}
\end{figure}\end{column}
\begin{column}{0.5\textwidth}
\vspace{2em}
	Key points on the Vityaz Trench: 
\begin{itemize}
            \item Vt has developed in a situation of complex tectonic settings.
            \item Tectonic plates convergence is marked by 2 opposite subduction zones, defined by the Vt (W) and the Tonga Trench (E)
            \item This results in an asymmetric opening of the North Fiji back-arc basin
            \item Subducting IAP and a subduction of the d’Entrecasteaux ridge $=>$ New Hebrides trench-arc system
            \item Southern NH Arc and NF Basin enlighten the geodynamic complexity and instability of the region
            \item Active volcano Mt. Yasur: located in Tanna Island of the Vanuatu arc
\end{itemize}
\end{column}
\end{columns}
\end{minipage}
\end{frame}

\section{NB}

\begin{frame}\frametitle{New Britain Trench}
\vspace{1em}
\begin{figure}[H]
	\centering
		\includegraphics[width=8cm]{Fig-NBT1.jpg}
	\caption{Topographic map and location of the New Britain and San Cristobal trenches}\label{fig:NBT1}
\end{figure}
\end{frame}

\begin{frame}\frametitle{New Britain Trench}
\vspace{1em}
\begin{figure}[H]
	\centering
		\includegraphics[width=8cm]{Fig-NBT6.jpg}
	\caption{Cross-section profiles of the New Britain and San Cristobal trenches}\label{fig:NBT6}
\end{figure}
\end{frame}

% ----------------------------------------------------------------------------
% *** START Philippine Trench <<<
% ----------------------------------------------------------------------------

\begin{frame}\frametitle{Philippine Trench}
\vspace{1em}
\begin{figure}[H]
	\centering
		\includegraphics[width=8.0cm]{Fig-Ph08.jpg}
	\caption{Plotting by Python (Matplotlib library): distribution of the bathymetric samples of the cross-section profiles, PhT}\label{fig:Ph08}
\end{figure}
%\vspace{1em}
\begin{exampleblock}{Data Processing by Python}
 \begin{dinglist}{46}\footnotesize
     \item Python can handle a huge amount of data, which it uses not only to process, but also to create analytical plots, such as assessment of bathymetric variability or analytics in geology. 
     \item Big data analysis in Python is based on a full range of algorithms used in statistics
     \item Statistical analysis: well developed stand-alone libraries (Matplotlib, SciPy, NumPy, Pandas, StatsModels)
     \item Beauty \& elegance of graphs: advanced design solutions for data visualization in Matplotlib \& Seaborn
     \item Scripting approach with more readable syntax
     \item Open-source
\end{dinglist}
\end{exampleblock}
\end{frame}

\begin{frame}\frametitle{Philippine Trench}
\begin{minipage}[0.4\textheight]{\textwidth}
\begin{columns}[T]
\begin{column}{0.5\textwidth}
\vspace{2em}
	\includegraphics[width=5.7cm]{Fig-Ph09.jpg}
 \begin{dinglist}{105}\footnotesize
     \item Example of data analysis by R: correlation matrix visualising \textcolor{red}{influence of impact factors} affecting the geomorphology of the profiles. Plotting: R.
\end{dinglist} 
\end{column}
\begin{column}{0.5\textwidth}
\vspace{2em}
	\includegraphics[width=5.7cm]{Fig-Ph10.jpg}
\begin{dinglist}{105}\footnotesize
     \item Analysing the interactions between the variables in big datasets by R presents a solution to \textcolor{red}{visualising heterogeneity} of topographic features formed under various geologic setting: ranking geomorphology of the cross-section profiles. 
\end{dinglist}
\end{column}
\end{columns}
\end{minipage}
\end{frame}

% ----------------------------------------------------------------------------
% *** END Philippine Trench <<<
% ----------------------------------------------------------------------------

\section{Mn}

\begin{frame}\frametitle{Manila Trench}

\begin{minipage}[0.4\textheight]{\textwidth}
\begin{columns}[T]
\begin{column}{0.5\textwidth}
\vspace{2em}
\begin{figure}[H]
	\centering
		\includegraphics[width=6.0cm]{Fig-M1.jpg}
	\caption{Topographic location of the Manila Trench and South China Sea}\label{fig:H1}
\end{figure}
\end{column}
\begin{column}{0.5\textwidth}
\vspace{2em}
\begin{itemize}
            \item 
\end{itemize}
\end{column}
\end{columns}
\end{minipage}
\end{frame}

 % ----------------------------------------------------------------------------
\section{R}

\begin{frame}\frametitle{Geologic mapping: Ryukyu Trench}
\begin{minipage}[0.4\textheight]{\textwidth}
\begin{columns}[T]
\begin{column}{0.5\textwidth}
\vspace{2em}
\begin{figure}[H]
	\centering
		\includegraphics[width=6.0cm]{Fig-R2.jpg}
	\caption{Geologic map of the Ryukyu Trench.}\label{fig:H1}
\end{figure}
\end{column}
\begin{column}{0.5\textwidth}
\vspace{2em}
\begin{itemize}
            \item 
\end{itemize}
\end{column}
\end{columns}
\end{minipage}
\end{frame}

 % ----------------------------------------------------------------------------
\section{Pl}

\begin{frame}\frametitle{Palau Trench}
\begin{minipage}[0.4\textheight]{\textwidth}
\begin{columns}[T]
\begin{column}{0.5\textwidth}
\vspace{2em}
\begin{figure}[H]
	\centering
		\includegraphics[width=6.0cm]{Fig-YPT1.jpg}
	\caption{Topographic location of the Yap and Palau Trenches}\label{fig:H1}
\end{figure}
\end{column}
\begin{column}{0.5\textwidth}
\vspace{2em}
	Key points on the Palau Trench: 
\begin{itemize}
            \item PT is located in the Philippine Sea, N from the Papua-New Guinea and E from the Philippines
            \item PT has greatest depth at 8,080 m (S), has straight geometric shape form
            \item PT and YT form a typical trench-arc-basin system which includes:
		\begin{itemize}
           		 \item deep trench, 
           		 \item island arcs, 
          		  \item back-arc basin (a Parece Vela Basin, located in the E side of the PSP).
         	 \end{itemize}
	\item Complex tectonic structure is reflected in notable heterogeneity of the crustal and upper mantle velocity structures beneath the PSP \cite{OdaSenna1994}. 
\end{itemize}
\end{column}
\end{columns}
\end{minipage}
\end{frame}

% ----------------------------------------------------------------------------
\section{Y}

\begin{frame}\frametitle{Yap Trench}
\begin{minipage}[0.4\textheight]{\textwidth}
\begin{columns}[T]
\begin{column}{0.6\textwidth}
\vspace{2em}
\begin{figure}[H]
	\centering
		\includegraphics[width=7.0cm]{Fig-YPT7.jpg}
	\caption{Automatically digitised profiles of Yap and Palau trenches}\label{fig:H1}
\end{figure}
\end{column}
\begin{column}{0.4\textwidth}
\vspace{2em}
	Key points on the Trench: 
\begin{itemize}
            \item YT is located in the Philippine Sea, N from the Papua-New Guinea and E from the Philippines.
            \item YT has the greatest depth at 8,292 m, is 700km long and 50km wide \cite{Satoetal1997}.
            \item YT stretches from the trench axis to Yap islands arc with a characteristic 'J' form \cite{Lee2004}. 
            \item YT is morphologically divided into the N and S parts with the the boundary between the two at $8^\circ26^\prime$N \cite{Fujiwaraetal2000}. 
            \item YT and PT are distinguished from the flat-bottomed Ayu Trough by their V shape. 
\end{itemize}
\end{column}
\end{columns}
\end{minipage}
\end{frame}

% ----------------------------------------------------------------------------
\section{MT}

\begin{frame}\frametitle{Mariana Trench}
\vspace{2em}
\begin{minipage}[0.4\textheight]{\textwidth}
\begin{columns}[T]
\begin{column}{0.4\textwidth}
 \begin{dinglist}{254}\small
     \item Large amount of Earth data brings a technical problem: how we can \textcolor{red}{extract necessary information} from the massive quantities (millions) of cells, pixels and variables?
     \item The answer to this question: effective means of statistical data analysis by \textcolor{red}{Python or R} (Julia, MATLAB/Octave) can help analyse, visualise and model data.
     \item \textcolor{red}{Automatisation of the massive digital data processing} for effective cartographic visualisation and statistical modelling. 
     \item As a response to such needs, \textcolor{red}{scripting cartography} and \textcolor{red}{statistical libraries} of the programming languages present effective and advanced approach in Earth data analysis and modelling. 
\end{dinglist}
\end{column}
\begin{column}{0.6\textwidth}
\begin{figure}[H]
	\centering
		\includegraphics[width=5.0cm]{Fig-MT20.jpg}
	\caption{Project Mindmap. Plotting: \LaTeX \space \cite{Lemenkova201993}}\label{fig:MT20}
\end{figure}	
\end{column}
\end{columns}
\end{minipage}
\end{frame}

 % ----------------------------------------------------------------------------
\begin{frame}\frametitle{Mariana Trench}
\vspace{2em}
\begin{minipage}[0.4\textheight]{\textwidth}
\begin{columns}[T]
\begin{column}{0.4\textwidth}
\begin{itemize}\footnotesize
	\item Various GMT techniques and modules have been used for rapid and accurate mapping
	\item Main GMT modules:
	\begin{dinglist}{224}\footnotesize
		\item  \textcolor{SlateBlue}{gmt makecpt} for making color palette
		\item  \textcolor{SlateBlue}{gmt grdimage} for generating raster image
		\item  \textcolor{SlateBlue}{gmt psscale} for adding legend for adding scale
		\item  \textcolor{SlateBlue}{gmt grdcontour} for adding isolines
		\item  \textcolor{SlateBlue}{gmt psbasemap} for adding grid
		\item  \textcolor{SlateBlue}{gmt psbasemap} for adding scale and directional rose
		\item  \textcolor{SlateBlue}{echo} UNIX utility for adding text labels
		\item  \textcolor{SlateBlue}{gmt pstext} for adding subtitle
		\item  \textcolor{SlateBlue}{gmt logo} for adding GMT logo
		\item  \textcolor{SlateBlue}{gmt psconvert} for convert PS to image (by GhostScript interpreter)
	 \end{dinglist}
\end{itemize}
\begin{itemize}\footnotesize
	\item As a result, plotting a map by GMT becomes a rapid procedure, fully controlled by the cartographer 
	\item Such automatisation by GMT facilitates mapping and decreases time of cartographic workflow
\end{itemize}
\end{column}
\begin{column}{0.6\textwidth}
\begin{figure}[H]
	\centering
		\includegraphics[width=6.0cm]{Fig-MT1.jpg}
	\caption{Topographic map and location of the Mariana Trench. \cite{Lemenkova201992}}\label{fig:MT1}
\end{figure}	
\end{column}
\end{columns}
\end{minipage}
\end{frame}

 % ----------------------------------------------------------------------------
\begin{frame}\frametitle{Mariana Trench}
\vspace{1em}
\begin{figure}[H]
	\centering
		\includegraphics[width=8cm]{Fig-MT11.jpg}
	\caption{Kernel Density Estimation (KDE) for the profiles bathymetry. \\
	Python libraries: Matplotlib, Seaborn, Pandas \cite{Rossumatal2018}. \cite{Lemenkova201905}}\label{fig:MT11}
\end{figure}
\begin{dinglist}{109}\footnotesize
     \item Example of Python based data analysis: Kernel Density Estimation (KDE) analysis showing facetted plots of the cross-section profiles. 
     \item It shows the range of density distribution in the elevations across the Mariana Trench
\end{dinglist} 
\end{frame}

 % ----------------------------------------------------------------------------
\begin{frame}\frametitle{Mariana Trench}
\vspace{1em}
\begin{figure}[H]
	\centering
		\includegraphics[width=10cm]{Fig-MT14.jpg}
	\caption{R based ridgeline plots by tectonic plates and bathymetric profiles, library \{ggplot\}. \cite{Lemenkova201902}}\label{fig:MT14}
\end{figure}

Key points on the ridgeline plot by R:
\begin{dinglist}{96}\footnotesize
	\item Analysing the interactions between the variables in big datasets by R presents a solution to \textcolor{red}{visualising heterogeneity} of topographic features formed under various geologic setting
	\item Ridgeline plot shows \textcolor{red}{data distribution} aimed to understand the variability of elevations as a \textcolor{red}{stream view}
	\item Ridgeline profiling shows a comparative visualisation of the \textcolor{red}{density distribution} of depths by profiles
	\item Ridgeline plot provides the insight into the statistical distribution of the \textcolor{red}{elevation ranges} in the Mariana Trench.
\end{dinglist}
\end{frame}

 % ----------------------------------------------------------------------------
\section{IBT}
\begin{frame}\frametitle{Izu-Bonin Trench}
\vspace{2em}
\begin{minipage}[0.4\textheight]{\textwidth}
\begin{columns}[T]
\begin{column}{0.4\textwidth}
Key points on the Izu-Bonin Trench (IBT): 
\begin{itemize}
            \item 
\end{itemize}
\end{column}
\begin{column}{0.5\textwidth}
\begin{figure}[H]
	\centering
		\includegraphics[width=6.0cm]{Fig-IBT1.jpg}
	\caption{Topographic map and location of the Izu-Bonin Trench}\label{fig:MAT1}
\end{figure}	
\end{column}
\end{columns}
\end{minipage}
\end{frame}

 % ----------------------------------------------------------------------------
\section{IBT}
\begin{frame}\frametitle{Izu-Bonin Trench}
\vspace{2em}
\begin{minipage}[0.4\textheight]{\textwidth}
\begin{columns}[T]
\begin{column}{0.4\textwidth}
\begin{figure}[H]
	\centering
		\includegraphics[width=5.5cm]{S04.jpg}
	\caption{Automatic digitizing of the IBT profiles: fragment of the GMT script}\label{fig:MAT1}
\end{figure}
\end{column}
\begin{column}{0.5\textwidth}
\begin{figure}[H]
	\centering
		\includegraphics[width=6.5cm]{Fig-IBT7.jpg}
	\caption{Cross-section profiles of the IBT}\label{fig:IBT7}
\end{figure}	
\begin{itemize}\scriptsize
            \item Automatic digitising by GMT enables visualising variations of geomorphology and model trends of the geomorphic shape curves of trenches
\end{itemize}
\end{column}
\end{columns}
\end{minipage}
\end{frame}

 % ----------------------------------------------------------------------------

\section{JT}

\begin{frame}\frametitle{Japan Trench}
\begin{minipage}[0.4\textheight]{\textwidth}
\begin{columns}[T]
\begin{column}{0.5\textwidth}
\vspace{2em}
	\includegraphics[width=6.0cm]{Fig-JT1.jpg}
The Japan Trench is located in W part of the Pacific Ocean, E off Japan. \\
Geologic and tectonic settings are strongly influenced by subduction of the PP beneath Japan and NAP along JT \cite{HueneLallemand1990}. 
\end{column}
\begin{column}{0.5\textwidth}
\vspace{2em}
	Key points on the Japan Trench (JT): 
\begin{itemize}
            \item JT region is a convergent zone associated with earthquakes \cite{Yoshii1979}. 
            \item JT belongs to circum-Pacific seismic belt 'Ring of Fire' $=>$ large earthquakes are in sedimentation area on upper forearc slopes \cite{Araietal2014}. 
            \item JT structure is a typical arc-trench system \cite{Ludwigetal1966}. 
            \item JT arc-trench system is a type of the mobile belts with bathymetric and topographic relief formed under the impact of tectonic plates subduction, active volcanism, repetitive earthquakes and orogenesis. 
            \item Geological complexity of the JT results in the highly seismicity \cite{HueneCulotta1989}; \cite{Linetal2016}; \cite{Bostonetal2017}.
\end{itemize}
\end{column}
\end{columns}
\end{minipage}
\end{frame}

\begin{frame}\frametitle{Japan Trench}
\begin{minipage}[0.4\textheight]{\textwidth}
\begin{columns}[T]
\begin{column}{0.5\textwidth}
\vspace{2em}
Key points on the cross-sectioned profiles visualized at two JT segments: 
	\includegraphics[width=6.5cm]{Fig-JT8.jpg}
rugged geomorphic structure of the median profiles and increased altitude in S-N direction on the N part
\end{column}
\begin{column}{0.45\textwidth}
\vspace{2em}	
\begin{itemize}
            \item N more steep slope on the continental slope part beneath the frontal wedge of the North American tectonic plate. 
            \item S shows more gentle slope on the Honshu Island. 
            \item 12 profiles were drawn for the southern segment and 14 profiles for the northern one.
            \item geological settings and seismicity of the area impact the trench roughness and general geomorphic shape
            \item W part of the trench profiles (Honshu Island slope) has steeper slopes comparing to the oceanward slopes for both segments
\end{itemize}
\end{column}
\end{columns}
\end{minipage}
\end{frame}

\section{Con}

\begin{frame}\frametitle{Research Innovation}
There are both theoretical and practical innovations of the presented research. 

\begin{itemize}
            \item The theoretical novelty lies in the comparative geomorphological mapping of the 20 Pacific hadal trenches that does not exists in the available literature. 
            \item The practical novelty consists in the developed methodology of the sequential technological chain of processes: GMT, QGIS plugins, Python \cite{Rossumatal2018}, Matlab/Octave, AWK and R.
            \item Cartographic novelty consists in the developed and presented algorithm of the cross-section profiles digitizing and geomorphic modelling of the  hadal trench modelling through the 'grdtrack' (main) GMT module. 
            \item Hadal trench present a complex system with highly interconnected factors affecting geomorphological structure, formation and development of the trench: slabs and tectonics plates, bathymetry, geographic location, geologic structure of the underlying basement and sediment thickness. 
            \item Therefore, comparative analysis of the trenches of the Pacific Ocean requires advanced methods of data analysis for operating with large data sets, structuring, organizing and managing thematic information in a GIS database, linking data and creating map overlays \cite{Lemenkova201540}, \cite{Lemenkova2016e}. 
            \item Integrating multi-source data supports verification of the data precision and control. The most important geodata include GEBCO \cite{Mayeretal2018}; \cite{Monahan2004}, SRTM \cite{Beckeretal2009}, ETOPO1 \cite{AmanteEakins2009}, Goole Earth \cite{Lemenkova201510}, CryoSat-2, Envisat, Jason-1 \cite{Sandwelletal2013}.
\end{itemize}

\end{frame}

\begin{frame}\frametitle{Research Significance and Justification}
The significance and justification of this works consist in the following:

\begin{itemize}
         \item Although ML has been significantly increased recently, using scripting and ML in cartography still remains lower comparing to the traditional GIS used in geosciences (e.g. \cite{Lemenkova2005b2}, \cite{RID09251802355041}, \cite{RID:0925180235504-2}, \cite{Lemenkova2014d}, \cite{Lemenkova2011c}, \cite{Lemenkova2005a}, \cite{Lemenkova2012j}, \cite{Lemenkova2012d}, \cite{Lemenkova2012g}). 
         \item Seafloor mapping includes multiple steps in technological process which may include: Hydrosweep DS sonar echo-sounding \cite{Lemenkova201569}, CARIS HIPS data processing \cite{Lemenkova2007f}, ArcGIS/ArcCatalog database management, GMT data processing \cite{Lemenkova201565}. 
	\item Accurate digitizing cross-section profiles using ArcGIS is slower comparing to the developed GMT-based automated digitizing. This is important for accurate geomorphological data analysis and crucial for better understanding of the seafloor landforms. 
	\item GMT-based mapping provides accurate visualization of the seafloor \cite{Lemenkova2007f}
	\item GMT-based mapping provides faster and preciser cartographic digitizing comparing to the traditional GIS. 
	\item The traditional handmade (e.g. via QGIS) cartographic digitizing is a tedious monotonous slow-speed routine, might be prone to minor or major errors. Efforts in developing methodology in automatization in bathymetric mapping exists \cite{Lemenkova2008b}. 
	\item The GMT based ML methodology (\cite{Wessel2019}, \cite{Wessel2018}) enabled to make further steps in automatization of the cartographic techniques. This significantly improved the speed, precision and quality of the data modelling.	
\end{itemize}

\end{frame}

\begin{frame}\frametitle{Conclusion}

\begin{itemize}
	\item Precise, correct and up-to-date information about the geomorphology hadal trenches in the Pacific Ocean is necessary understanding marine geology, tectonics, seismicity, processes of sedimentation and geodynamic evolution of the seafloor \cite{Lemenkova2006e}, \cite{Lemenkova2006a}, \cite{Lemenkova2006b}, \cite{Lemenkova2006c}.
	\item Submarine geomorphological and bathymetric mapping plays a critical role in analysis of the geological structures of the seafloor, marine benthic habitats, navigation, geological drilling, modelling marine environment and other aspects of geosciences.	
	\item Current studies contributed to the methodological testing and technical application of the advanced algorithms for seafloor modelling and mapping and to the geomorphological modelling. 	
	 \item Tested, presented and explained functionality of the several GMT modules enables to do automated digitizing of the orthogonal profiles crossing trenches in the perpendicular direction. Through this modelling, the shape of the landforms and steepness gradient of the trenches were visualized, compared and statistically analyzed. 
\end{itemize}

\end{frame}

% ----------------------------------------------------------------------------
\section{Thx}
\begin{frame}\frametitle{Thanks}
\begin{tikzpicture}%
\node[anchor=south west, inner sep=0] (X) at (-0,0){\includegraphics[width=11.5cm]{Fig-Final.jpg}};%
\begin{scope}[x={(X.south east)},y={(X.north west)}]%
%% This makes all measurements with no units into fractions of the width
%% and height of the graphic contained in the node.
\fill[fill = white,fill opacity=0.6] ($(X.south west) + (0,1in)$) rectangle ($(X.north east) - (0,1in)$);%
\node[text width=8cm] (Z) at (0.5,0.5) {%
  \sffamily\centering\large \textcolor{NavyBlue}{Thank you for attention !}\\[3pt]
  \normalsize \textcolor{NavyBlue}{Questions ?}\par%
};
\end{scope}%
\end{tikzpicture}

\end{frame}

% ----------------------------------------------------------------------------
% *** START: Final Slide <<<
% ----------------------------------------------------------------------------

%\section{End}
%\begin{frame}\frametitle{Thank you for attention!}

%\begin{figure}[H]
%	\centering
%		\includegraphics[width=7cm]{Wordcloud.jpg}
%\end{figure}	

%\end{frame}

%\section{Thx}
%\begin{frame}\frametitle{Acknowledgements}

%I gratefully acknowledge the funding of this research from the China Scholarship Council (CSC),
%State Ocean Administration (SOA), Marine Scholarship of China, People's Republic of China 
%(P. R. C.), Beijing, Grant \#2016SOA002, 2016-2020.

%\end{frame}
 
% ----------------------------------------------------------------------------
% *** END: Final Slide >>>
% ----------------------------------------------------------------------------
%%%%%%%%%%% Bibliography %%%%%%%
\section{Bib}
	\markright{Bibliography}
	\printbibliography[heading=none]
	
\end{document}

%Changing the font size locally (from biggest to smallest):	
%\Huge
%\huge
%\LARGE
%\Large
%\large
%\normalsize (default)
%\small
%\footnotesize
%\scriptsize
%\tiny